\section{Posisi dan Kebaruan Penelitian}

Tabel \ref{tab:posisi-penelitian} menunjukkan dua jalur penelitian yang berkembang secara terpisah. 
Jalur pertama mencakup \textit{ABAC policy mining} berbasis evolusioner atau metaheuristik \parencite{gaspar-cunha_evolutionary_2015,shang_abac_2024,li_attribute_2026} serta non-evolusioner \parencite{guo_--fly_2025}, yang belum menjadikan keterbatasan memori perangkat \textit{edge-IoT} sebagai pertimbangan utama desain. 
Jalur kedua mencakup perkembangan EDA dan varian hemat memori, seperti \textit{compact-GA} dan \textit{compact-PSO} \parencite{khalfi_metaheuristics_2023}, serta pemodelan dependensi melalui BOA, iBOA, dan pohon Chow-Liu \parencite{pelikan_iboa_2008,chow_approximating_1968}. 
Berdasarkan penelusuran literatur, pendekatan tersebut belum diterapkan pada \textit{ABAC policy mining}. 
CoDA-EDA diposisikan pada titik temu kedua jalur tersebut.

CoDA-EDA dibangun berdasarkan tiga prinsip. 
Pertama, \textit{compact}, yaitu menggunakan pembaruan pemenang-pecundang pada \textit{Probability Vector} (PV) tanpa menyimpan populasi eksplisit \parencite{khalfi_metaheuristics_2023}, sehingga sesuai dengan keterbatasan sumber daya perangkat \textit{edge-IoT} \parencite{canavese_security_2024}. 
Kedua, \textit{dependency-aware}, yaitu mempertahankan pemodelan dependensi antaratribut dengan struktur yang dibatasi, sebagai kompromi antara asumsi independensi pada algoritma \textit{compact} dan pemodelan multivariat yang lebih kompleks pada BOA/iBOA \parencite{pelikan_iboa_2008}. 
Ketiga, EDA, yaitu menempatkan kedua mekanisme tersebut dalam kerangka seleksi, estimasi, dan \textit{sampling} berbasis distribusi probabilistik.

Dinamika lingkungan IoT \parencite{ragothaman_access_2023} juga menuntut kebijakan tetap adaptif terhadap perubahan data. 
Untuk menghindari biaya pembelajaran ulang struktur pada setiap siklus, CoDA-EDA memisahkan pembaruan parameter dan struktur dependensi. 
Parameter probabilistik diperbarui setiap siklus, sedangkan struktur hanya dipelajari ulang ketika terdeteksi pergeseran distribusi yang signifikan. 
Mekanisme ini berbeda dari iBOA yang memperbarui struktur dan parameter secara inkremental dalam satu kerangka pembaruan \parencite{pelikan_iboa_2008}, serta menjadi dasar evaluasi efisiensi komputasi dan kemampuan \textit{warm-start} ketika terjadi perubahan data.

Deteksi perubahan tersebut memerlukan ukuran perbedaan distribusi antarjendela. 
Penelitian ini menggunakan \textit{Jensen–Shannon divergence} (JSD)\label{first:jsd} untuk mengukur besarnya pergeseran distribusi. 
Untuk dua distribusi probabilitas $P$ dan $Q$ dengan bobot yang sama, JSD dirumuskan sebagai:

\begin{equation}
    \operatorname{JSD}(P \parallel Q)
    =
    \frac{1}{2}
    \operatorname{KL}(P \parallel M)
    +
    \frac{1}{2}
    \operatorname{KL}(Q \parallel M),
    \qquad
    M = \frac{1}{2}(P+Q),
    \label{eq:jsd-drift}
\end{equation}

dengan $KL$ menyatakan \textit{Kullback–Leibler divergence}. 
JSD dipilih karena memiliki sifat yang sesuai untuk membandingkan distribusi diskret. 
Berbeda dari \textit{KL divergence}, kelas divergensi berbasis entropi Shannon yang mencakup JSD tidak mensyaratkan kontinuitas absolut, sehingga tetap terdefinisi ketika suatu nilai atribut memiliki probabilitas nol pada salah satu distribusi \parencite{lin1991divergence}. 
Dengan logaritma berbasis dua dan bobot yang sama, JSD juga bersifat non-negatif, simetris, dan terbatas pada rentang 0–1; nilai 0 menunjukkan distribusi identik, sedangkan nilai yang lebih besar menunjukkan perbedaan distribusi yang lebih tinggi \parencite{lin1991divergence}. 
Sifat berbatas dan simetris tersebut memudahkan penggunaan JSD sebagai dasar penetapan ambang perubahan distribusi. 
\textcite{lin1991divergence} juga menunjukkan hubungan JSD dengan batas probabilitas kesalahan klasifikasi Bayes, yang memperkuat interpretasinya sebagai ukuran keterbedaan informasi antar-distribusi.

Penggunaan JSD untuk mendeteksi perubahan distribusi telah diterapkan pada concept \textit{drift}. 
\textcite{yang2022concept} mengusulkan CDJD (\textit{Concept Drift Detection based on Jensen–Shannon Divergence}) dengan membandingkan distribusi fitur antara jendela lama dan jendela berjalan. 
\textit{Concept drift} dinyatakan ketika JSD melampaui ambang tertentu, sedangkan kemiripan dengan konsep terdahulu dapat digunakan untuk memanfaatkan kembali model yang tersedia. 
Hasil eksperimen menunjukkan bahwa pendekatan tersebut mampu mendeteksi perubahan distribusi dan mempertahankan kinerja klasifikasi pada aliran data yang mengalami \textit{concept drift} \parencite{yang2022concept}.

CoDA-EDA mengadopsi prinsip tersebut untuk tujuan yang berbeda. 
JSD digunakan untuk mengukur pergeseran distribusi pasangan atribut–nilai antarjendela log akses. 
Struktur dependensi Chow–Liu dipelajari ulang hanya ketika JSD melampaui ambang $\tau$\label{first:symbol-tau}, sedangkan parameter probabilistik tetap diperbarui pada setiap siklus. 
Dengan demikian, JSD berfungsi sebagai pemicu yang membedakan perubahan distribusi yang memerlukan pembaruan struktur dari variasi yang masih dapat ditangani melalui pembaruan parameter. 
Mekanisme ini mengurangi pembelajaran ulang struktur yang tidak diperlukan pada perangkat \textit{edge} dengan sumber daya terbatas.

Dibandingkan penelitian pada Tabel \ref{tab:posisi-penelitian}, CoDA-EDA berbeda dari pendekatan berbasis populasi \textcite{gaspar-cunha_evolutionary_2015}, \textcite{shang_abac_2024}, dan \textcite{li_attribute_2026} karena menggunakan pencarian probabilistik tanpa menyimpan populasi eksplisit. 
Perbedaan ini menjadi dasar evaluasi jejak memori terhadap BOA, iBOA, ACO, dan UBK. 
CoDA-EDA juga berbeda dari \textcite{guo_--fly_2025}, yang menggunakan optimasi submodular deterministik/\textit{greedy} dengan fokus pada usabilitas kebijakan bagi manusia dan LLM. 
Sebaliknya, CoDA-EDA dirancang sebagai pencarian probabilistik berbasis model dengan fokus pada efisiensi sumber daya perangkat \textit{edge}.

Diferensiasi CoDA-EDA terhadap EDA terdahulu tidak terletak pada satu komponen, tetapi pada kombinasi beberapa keputusan rancangan. 
Tabel \ref{tab:perbedaan-coda-eda-long} membandingkan CoDA-EDA dengan keluarga EDA terdekat berdasarkan empat dimensi, yaitu: (1) model dependensi antaratribut, (2) penyimpanan populasi dan kebutuhan memori terhadap ukuran populasi virtual (NP), (3) mekanisme pembaruan struktur model, dan (4) domain penerapan.

\begin{longtable}{|>{\raggedright\arraybackslash}p{2cm}|>{\raggedright\arraybackslash}p{3cm}|>{\raggedright\arraybackslash}p{3cm}|>{\raggedright\arraybackslash}p{3cm}|>{\centering\arraybackslash}p{2cm}|}
\caption{Perbedaan CoDA-EDA terhadap Algoritma EDA yang Sudah Ada}
\label{tab:perbedaan-coda-eda-long} \\
\hline
\textbf{Algoritma EDA} & \textbf{Model dependensi antar-atribut} & \textbf{Populasi \& memori terhadap NP} & \textbf{Mekanisme pembaruan struktur} & \textbf{Diterapkan pada ABAC?} \\
\hline
\endfirsthead

\multicolumn{5}{c}{{\bfseries \tablename\ \thetable{} -- lanjutan dari halaman sebelumnya}} \\
\hline
\textbf{Algoritma EDA} & \textbf{Model dependensi antar-atribut} & \textbf{Populasi \& memori terhadap NP} & \textbf{Mekanisme pembaruan struktur} & \textbf{Diterapkan pada ABAC?} \\
\hline
\endhead

\hline
\multicolumn{5}{r}{{Lanjut ke halaman berikutnya}} \\
\endfoot

\hline
\endlastfoot

% Data
\textit{compact-GA / compact-PSO} & Tidak ada --- univariat, independensi penuh & Tanpa populasi eksplisit; \textit{Probability Vector} O(D), memori datar terhadap NP & Tidak ada struktur untuk diperbarui & Belum \\
\hline
BOA & Jaringan Bayes tanpa batasan orde (multivariat penuh) & Populasi eksplisit; memori tumbuh terhadap NP & Membangun ulang jaringan Bayes tiap generasi (mahal secara komputasi) & Belum \\
\hline
iBOA & Pohon dependensi (Chow-Liu), inkremental & Statistik/populasi inkremental; tidak menanggung transien populasi penuh & Struktur dan parameter diperbarui pada skema inkremental yang sama, tiap langkah & Belum \\
\hline
CoDA-EDA (diusulkan) & Pohon dependensi Chow-Liu dibatasi order-1 (terbatas, bukan tanpa batas) & Tanpa populasi eksplisit; \textit{Probability Vector} + pohon O(D), memori datar terhadap NP & Struktur dipisah dari parameter: struktur jarang (dipicu \textit{drift}), parameter tiap siklus & Ya --- penelitian ini (Raspberry Pi 5) \\
\hline

\end{longtable}

Tabel \ref{tab:perbedaan-coda-eda-long} memperjelas posisi CoDA-EDA pada setiap dimensi pembanding. 
Dari sisi dependensi, CoDA-EDA berada di antara pendekatan \textit{compact} yang mengasumsikan independensi penuh dan BOA yang memodelkan dependensi multivariat tanpa batasan orde, dengan menggunakan pohon Chow–Liu berorde satu. 
Dari sisi memori, CoDA-EDA mempertahankan kebutuhan yang tidak bergantung pada ukuran populasi virtual (NP) sebagaimana keluarga \textit{compact}, tetapi tetap memodelkan dependensi antarvariabel. 
Perbedaan utama terhadap iBOA terletak pada mekanisme pembaruan model: parameter probabilistik diperbarui setiap siklus, sedangkan struktur dependensi hanya dipelajari ulang ketika terdeteksi \textit{drift}. 
Kombinasi karakteristik tersebut menjadi dasar kontribusi metodologis CoDA-EDA.

Kontribusi tersebut terletak pada penggabungan tiga mekanisme yang sebelumnya digunakan secara terpisah dalam literatur EDA, yaitu pembaruan tanpa populasi eksplisit ala \textit{compact}, pemodelan dependensi berorde terbatas melalui pohon Chow–Liu, dan pemisahan frekuensi pembaruan struktur serta parameter. 
Berdasarkan sintesis tersebut, penelitian ini mengusulkan tiga pilar kebaruan.

Kebaruan pertama adalah penerapan EDA yang ringan dan peka terhadap dependensi pada \textit{ABAC policy mining}, yang berdasarkan penelusuran literatur belum ditemukan pada domain tersebut. 
Kebaruan kedua adalah penggabungan pemodelan dependensi orde pertama menggunakan pohon Chow–Liu dengan mekanisme pembaruan tanpa populasi eksplisit. 
Kebaruan ketiga adalah pemisahan pembaruan parameter yang dilakukan setiap siklus dari pembelajaran ulang struktur yang hanya dipicu ketika terjadi pergeseran distribusi.

Rancangan tersebut selanjutnya dievaluasi pada perangkat \textit{edge-IoT} nyata, yaitu Raspberry Pi 5, untuk menguji apakah efisiensi memori dan komputasi dapat dicapai tanpa mengorbankan kualitas kebijakan secara signifikan dibandingkan metode pembanding. 
Evaluasi ini melengkapi studi pada Tabel \ref{tab:posisi-penelitian} yang tidak menempatkan keterbatasan sumber daya perangkat \textit{edge} sebagai fokus utama rancangan algoritma.